% GENERATED FILE. Edit canonical lessons, then run npm run generate:books.
% canonical-source-hash: fnv1a64:abb05010f4dff1c4
% canonical-lessons: IT-C453-ferramenta, IT-C453-merceria, IT-C453-profumeria, IT-C453-erboristeria, IT-C453-gioielleria

\chapter{Piccoli negozi}
\label{ch:it-piccoli-negozi}
\hlchaptermodality{453}

\begin{chapteropening}
\textbf{By the end of this chapter:} \emph{I can find specialist shops: the hardware shop, the haberdasher, the perfume shop, the herbalist and the jeweller.}

Complete the last lesson of chapter 453: I can find specialist shops: the hardware shop, the haberdasher, the perfume shop, the herbalist and the jeweller.
\end{chapteropening}

\section[la ferramenta]{la ferramenta — the hardware shop; hardware}

\label{lesson:IT-C453-ferramenta}



\begin{quote}
Before the new one: say the Italian for the newsagent, then the Italian for the stationer's.
\end{quote}

\subsection*{You'll want to know: la ferramenta}
\textbf{la ferramenta} — \textquotedblleft{}the hardware shop; hardware\textquotedblright{}.

\begin{quote}
In ferramenta ho trovato le viti giuste.
\end{quote}

\begin{quote}
\emph{At the hardware shop I found the right screws.}
\end{quote}

\emph{Il negozio di ferramenta} in full; \emph{la ferramenta} is also the goods.

\begin{cousinweb}
\emph{ferramenta} is inherited from Latin \emph{ferrāmentum} 'iron tool' (EtymDB), from \emph{ferrum} 'iron'.
\end{cousinweb}

\begin{grammarlens}[title={What you've built}]
One more word for telling, explaining and describing at B1.
\end{grammarlens}

\subsection*{Your turn}
\begin{itemize}
\raggedright
  \item \emph{Say it:} \emph{la ferramenta}
  \item \emph{Say it:} \emph{la ferramenta}, once more
  \item \emph{Say it:} a sentence of your own with \emph{la ferramenta}
  \item \emph{Recall:} say the Italian for the newsagent, then the Italian for the stationer's, then say \emph{la ferramenta} again
\end{itemize}

\begin{tcolorbox}[breakable,colback=teal!4,colframe=teal!35!black,title={Before you move on}]
What does la ferramenta mean? (\textquotedblleft{}The hardware shop; hardware\textquotedblright{}.) Say it once more.
\end{tcolorbox}
\section[la merceria]{la merceria — the haberdashery}

\label{lesson:IT-C453-merceria}



\begin{quote}
Before the new one: say the Italian for the stationer's, then the Italian for the hardware shop.
\end{quote}

\subsection*{You'll want to know: la merceria}
\textbf{la merceria} — \textquotedblleft{}the haberdashery\textquotedblright{}.

\begin{quote}
In merceria vendono bottoni e filo.
\end{quote}

\begin{quote}
\emph{At the haberdasher's they sell buttons and thread.}
\end{quote}

Buttons, thread, ribbons and zips.

\begin{cousinweb}
\emph{merceria} is built on \emph{merce} 'goods', which is inherited from Latin \emph{merx, mercis} 'goods' (Lewis \& Short).
\end{cousinweb}

\begin{grammarlens}[title={What you've built}]
One more word for telling, explaining and describing at B1.
\end{grammarlens}

\subsection*{Your turn}
\begin{itemize}
\raggedright
  \item \emph{Say it:} \emph{la merceria}
  \item \emph{Say it:} \emph{la merceria}, once more
  \item \emph{Say it:} a sentence of your own with \emph{la merceria}
  \item \emph{Recall:} say the Italian for the stationer's, then the Italian for the hardware shop, then say \emph{la merceria} again
\end{itemize}

\begin{tcolorbox}[breakable,colback=teal!4,colframe=teal!35!black,title={Before you move on}]
What does la merceria mean? (\textquotedblleft{}The haberdashery\textquotedblright{}.) Say it once more.
\end{tcolorbox}
\section[la profumeria]{la profumeria — the perfume shop, the perfumery}

\label{lesson:IT-C453-profumeria}



\begin{quote}
Before the new one: say the Italian for the hardware shop, then the Italian for the haberdashery.
\end{quote}

\subsection*{You'll want to know: la profumeria}
\textbf{la profumeria} — \textquotedblleft{}the perfume shop, the perfumery\textquotedblright{}.

\begin{quote}
In profumeria ho comprato una crema.
\end{quote}

\begin{quote}
\emph{At the perfume shop I bought a cream.}
\end{quote}

Make-up and soaps too.

\begin{cousinweb}
\emph{profumeria} is built on \emph{profumo} 'scent', from \emph{per} 'through' and \emph{fumo} 'smoke', which is inherited from Latin \emph{fūmus} (Lewis \& Short): scent carried through smoke.
\end{cousinweb}

\begin{grammarlens}[title={What you've built}]
One more word for telling, explaining and describing at B1.
\end{grammarlens}

\subsection*{Your turn}
\begin{itemize}
\raggedright
  \item \emph{Say it:} \emph{la profumeria}
  \item \emph{Say it:} \emph{la profumeria}, once more
  \item \emph{Say it:} a sentence of your own with \emph{la profumeria}
  \item \emph{Recall:} say the Italian for the hardware shop, then the Italian for the haberdashery, then say \emph{la profumeria} again
\end{itemize}

\begin{tcolorbox}[breakable,colback=teal!4,colframe=teal!35!black,title={Before you move on}]
What does la profumeria mean? (\textquotedblleft{}The perfume shop, the perfumery\textquotedblright{}.) Say it once more.
\end{tcolorbox}
\section[l'erboristeria]{l'erboristeria — the herbalist's}

\label{lesson:IT-C453-erboristeria}



\begin{quote}
Before the new one: say the Italian for the haberdashery, then the Italian for the perfume shop.
\end{quote}

\subsection*{You'll want to know: l'erboristeria}
\textbf{l'erboristeria} — \textquotedblleft{}the herbalist's\textquotedblright{}.

\begin{quote}
In erboristeria vendono tisane e piante.
\end{quote}

\begin{quote}
\emph{At the herbalist's they sell herbal teas and plants.}
\end{quote}

\emph{L'erborista} — the herbalist.

\begin{cousinweb}
\emph{erboristeria} is built on \emph{erba} 'herb, grass', which is inherited from Latin \emph{herba} (EtymDB).
\end{cousinweb}

\begin{grammarlens}[title={What you've built}]
One more word for telling, explaining and describing at B1.
\end{grammarlens}

\subsection*{Your turn}
\begin{itemize}
\raggedright
  \item \emph{Say it:} \emph{l'erboristeria}
  \item \emph{Say it:} \emph{l'erboristeria}, once more
  \item \emph{Say it:} a sentence of your own with \emph{l'erboristeria}
  \item \emph{Recall:} say the Italian for the haberdashery, then the Italian for the perfume shop, then say \emph{l'erboristeria} again
\end{itemize}

\begin{tcolorbox}[breakable,colback=teal!4,colframe=teal!35!black,title={Before you move on}]
What does l'erboristeria mean? (\textquotedblleft{}The herbalist's\textquotedblright{}.) Say it once more.
\end{tcolorbox}
\section[la gioielleria]{la gioielleria — the jeweller's}

\label{lesson:IT-C453-gioielleria}



\begin{quote}
Before the new one: say the Italian for the perfume shop, then the Italian for the herbalist's.
\end{quote}

\subsection*{You'll want to know: la gioielleria}
\textbf{la gioielleria} — \textquotedblleft{}the jeweller's\textquotedblright{}.

\begin{quote}
In gioielleria guardiamo gli anelli.
\end{quote}

\begin{quote}
\emph{At the jeweller's we gaze at the rings.}
\end{quote}

\emph{Il gioiello} — a jewel.

\begin{cousinweb}
\emph{gioielleria} is built on \emph{gioiello} 'jewel', which was borrowed from Old French \emph{joel}, usually traced to Latin \emph{iocus} 'game, joke' (Lewis \& Short).
\end{cousinweb}

\begin{grammarlens}[title={What you've built}]
One more word for telling, explaining and describing at B1.
\end{grammarlens}

\subsection*{Your turn}
\begin{itemize}
\raggedright
  \item \emph{Say it:} \emph{la gioielleria}
  \item \emph{Say it:} \emph{la gioielleria}, once more
  \item \emph{Say it:} a sentence of your own with \emph{la gioielleria}
  \item \emph{Recall:} say the Italian for the perfume shop, then the Italian for the herbalist's, then say \emph{la gioielleria} again
\end{itemize}

\begin{tcolorbox}[breakable,colback=teal!4,colframe=teal!35!black,title={Before you move on}]
What does la gioielleria mean? (\textquotedblleft{}The jeweller's\textquotedblright{}.) Say it once more.
\end{tcolorbox}
